% comp s07 (W4, phase 2): the causal core; log misc/work-orders/comp-phase2-W4.applied.md
\section{Posteriors of effects from simulated pairs}\label{gbc:sec:efficiency}\label{comp:sec:core}

The network does not know what an average treatment effect is. It learns the law of whatever the simulator puts in the target, given whatever it puts in the summary, so the causal content of the analysis is decided before training starts, by six decisions\index{generative Bayesian computation!posteriors of effects|textbf}.

\begin{enumerate}[leftmargin=*]
\item \emph{Choose the estimand.} The simulator draws both potential outcomes, so it can compute the model based effect (the average of $\tau(X_i;\theta)$ over the observed covariates), the sample effect (the average of the simulated $Y_i(1)-Y_i(0)$, whose observed halves are part of the simulated data) or the population effect. The posterior of the sample effect depends on the joint law assumed for the pair and need not be wider than that of the model based effect (Section~\ref{sec:imputation}).
\item \emph{Train on it.} Train on $\psi(\theta)$ for the model based or population effect, or the average simulated $Y_i(1)-Y_i(0)$ for the sample effect, to learn its marginal posterior directly. Alternatively, push joint draws of $\theta$ through $\psi$ using the triangular map (Section~\ref{comp:sec:noise}). Separate marginal networks lose dependence and give the wrong variance for functionals of correlated coordinates.
\item \emph{Learn the effect as a function.} Give the network $x$ as an input and draw $x$ at random in the simulator, or train at a grid. A joint statement about the curve needs joint draws of its values at the grid points, from a triangular map over the points (Section~\ref{comp:sec:noise}) or from joint draws of $\theta$; independent levels at each point lose the dependence.
\item \emph{Put the estimator in the summary.} With an AIPW or H\'ajek estimate in $S$ (Chapter~\ref{ch:ipw}) the exact summary posterior is asymptotically at least as concentrated as the estimator (Section~\ref{comp:sec:summaries}). Cross fit its nuisances inside the simulator, since residuals on the training simulations are too small; in Chapter~\ref{ch:blp} cross fitting raised the average coverage of nominal 90 percent intervals from 0.85 to 0.94.
\item \emph{Simulate the confounding we fear.} Calibration on simulated data certifies the computation inside the simulator, not the simulator against the world. The prior predictive must contain confounding of the strength and form the application plausibly has, under the coupled prior of Chapter~\ref{ch:horseshoe}.
\item \emph{Check.} Run the checks of Section~\ref{comp:sec:checks}, against an exact reference where there is one.
\end{enumerate}

\paragraph{What holds.} In a correctly specified regular parametric model, with a prior positive and continuous at the truth and summaries that are sufficient or contain an efficient estimator, the Bernstein von Mises theorem gives a posterior that is asymptotically normal, centred at an efficient estimator, with the inverse Fisher information as variance. By the delta method any smooth functional, such as an average effect $\tau_Q(\theta)$, then has a posterior efficient for that model, and the network inherits this up to its training error, which no theorem bounds.

\paragraph{What does not hold automatically.} Semiparametric efficiency, the bound of Hahn (1998) for the nonparametric model of the regression and the propensity score, attained by the estimator of Hirano, Imbens and Ridder (2003), is another matter. A prior that ignores the propensity score can concentrate around the right regression function and still give an average effect biased at the $\sqrt N$ scale, when the regression is rough and the propensity informative (Ray and van der Vaart 2020), and a shrinkage prior can produce regularization induced confounding (Hahn and coauthors 2018). Generative computation changes how the posterior is computed, not which posterior it is: with such a prior every map returns the same wrong posterior. The repair is in the prior or in the estimator: Chapter~\ref{ch:horseshoe} couples the prior to the treatment model, with the limits that chapter reports for large $p$, and Section~\ref{sec:baipw} puts the horseshoe inside the doubly robust score of DML.

\subsection{A worked example with an exact posterior}

The test model is the stylized supplement trial of Chapter~\ref{ch:cdt}, after the Lanarkshire milk experiment (Section~\ref{sec:milk}): $n=200$ children with a baseline measure of undernourishment $Z_i\sim N(0,1)$ and a supplement $W_i$ given with probability $\expit(1+0.8Z_i)$. Covariates and assignments are drawn once and held fixed: 140 treated and 60 control children, with mean $Z$ of 0.26 and $-0.72$ respectively. The outcome is $Y_i=\mu_w+\alpha_wZ_i+\varepsilon_i$ for $w=W_i$, with $\varepsilon_i\sim N(0,1)$ and the variance known, and $\theta=(\mu_0,\alpha_0,\mu_1,\alpha_1)$ has the prior $N(0,4I)$. The targets are the model based effect $\psi=(\mu_1-\mu_0)+(\alpha_1-\alpha_0)\bar Z$, with $\bar Z$ the mean of the $Z_i$, and the conditional effect $\tau(z)=(\mu_1-\mu_0)+(\alpha_1-\alpha_0)z$.

The four within arm least squares coefficients $S$ are sufficient, so every network error is amortization error (Section~\ref{comp:sec:checks}). The arms have independent normal posteriors: with design $B_w$ (rows $(1,Z_i)$) and least squares coefficients $b_w$,
\[
(\mu_w,\alpha_w)^\top\mid S \sim N\bigl(V_wB_w^\top B_wb_w,\ V_w\bigr),\qquad V_w=\bigl(B_w^\top B_w+\tfrac14 I_2\bigr)^{-1}.
\]
The covariance depends only on the fixed design, so the exact posterior sd of $\psi$, the square root of $(1,\bar Z)(V_0+V_1)(1,\bar Z)^\top$, is 0.179 at every data set, against a prior sd of 2.83.

The ATE network is the quantile network of Section~\ref{comp:sec:network} with $\psi(\theta)$ as the target, trained with the check loss on 100,000 simulated pairs of $\theta$ and $S$ for 40 passes, about 26 seconds on two threads of an Intel server node. The observed data set is one draw of $\theta$ from the prior and the data simulated at it; the checks use 500 fresh studies with 100 draws each and 20,000 fresh pairs. Each network was trained once, so the observed data set results are for one data set and one seed.

At the observed data set the exact posterior of $\psi$ has mean 0.079 and sd 0.179, and 20,000 draws from the network have mean 0.067 and sd 0.185. Their Wasserstein-1 distance is 0.015, or 0.085 of the posterior sd (Figure~\ref{comp:fig:ate}a and c).

\begin{figure}[ht]
\centering
\includegraphics[width=\linewidth,trim=0bp 185bp 0bp 0bp,clip]{gbc_ate.pdf}
\caption{The ATE network of the worked example. (a) Draws of $\psi$ at the observed data set, the exact posterior density and the true value (dashed line). (b) Rank histograms over 500 studies for the ATE network, the exact posterior and the prior only network (trained with the summary replaced by a constant); the line marks the expected count of 50. (c) Quantile functions of the same three at the observed data set. (d) Wasserstein-1 error over the exact posterior sd against the Mahalanobis distance of $S$ from the center of the training summaries, for the 500 studies (the 50 farthest as squares) and along a ray leaving the center in the direction in which the exact mean of $\psi$ moves fastest; the dotted line is the largest distance among the training summaries.}
\label{comp:fig:ate}
\end{figure}

The rank check over 500 studies (Table~\ref{comp:tab:ate}, Figure~\ref{comp:fig:ate}b) cannot tell the ATE network from the exact posterior. Nor can it tell either from a prior only network, which is trained with the summary replaced by a constant and has intervals about 16 times as wide. The held out loss can: it puts the prior only network at 0.798 against 0.050 for the exact quantile function, and resolves the ATE network's excess, 0.00024, at 8.9 standard errors.

\begin{table}[ht]
\centering
\caption{Checks of three posteriors for $\psi$. Ranks: 500 studies with 100 draws each, in ten bins. Intervals run between the 5th and 95th percentiles of the draws (Weibull plotting positions). Held out loss: the check loss on 20,000 fresh pairs with eight random levels each, for the exact quantile function and for the raw network function (0.05027 for the ATE network after sorting).}
\label{comp:tab:ate}
\footnotesize\setlength{\tabcolsep}{4pt}
\begin{tabular}{L{1.2in}>{\raggedleft\arraybackslash}p{0.8in}>{\raggedleft\arraybackslash}p{0.55in}>{\raggedleft\arraybackslash}p{0.55in}>{\raggedleft\arraybackslash}p{0.6in}}
\toprule
 & Rank $\chi^2$ ($p$) & Coverage of 90\% & Width of 90\% & Held out loss \\
\midrule
Exact posterior & 10.12 (0.341) & 0.876 & 0.603 & 0.05012 \\
ATE network & 8.12 (0.522) & 0.870 & 0.601 & 0.05037 \\
Prior only network & 8.36 (0.498) & 0.882 & 9.514 & 0.79757 \\
\bottomrule
\end{tabular}
\end{table}

The 50 studies with summaries farthest from the training center (Mahalanobis distance) have a flat rank histogram and a Wasserstein-1 error of 0.110 of the exact sd, against 0.061 for the other 450, although all 500 lie inside the training range (largest distance 4.59, against 5.61 among the training summaries). Beyond that range, along the ray of Figure~\ref{comp:fig:ate}d, the error passes the posterior sd (1.19 of it at distance 10).

The CATE network, with $z$ as a further input drawn from $N(0,1)$ for each simulated pair, is less accurate than the ATE network at this data set: its Wasserstein-1 errors are 0.114, 0.531 and 0.128 of the exact sd at $z=-1.5$, 0 and 1.5 (Figure~\ref{comp:fig:cate}e).

The 0.531 is mostly one unlucky fit. We trained ten networks of each kind, differing in initial weights and minibatch order, at each of 10{,}000, 40{,}000, 100{,}000 and 400{,}000 simulated pairs (each training set a prefix of the next and the same for every seed, with the architecture and 40 passes above), and compared the sorted network quantiles on 999 levels from 0.001 to 0.999 with the exact posterior quantiles on 2{,}000 fresh studies, by the Wasserstein-1 error over the exact sd. Each network gives a median and a 90th percentile of this error over the studies and a coverage, and we report their means over the ten seeds.

At 100{,}000 pairs the CATE network at $z=0$ has median error 0.15 (0.14 at $z=-1.5$, 0.11 at $z=1.5$), 90th percentile 0.33 and coverage of 90 percent intervals 0.91. At the observed data set its error averages 0.15 over the ten seeds, with standard deviation 0.13 and largest value 0.49, so our single network, trained on a different draw of simulations, lies above all ten. The error falls roughly as $J^{-1/2}$, from a median of 0.60 at 10{,}000 pairs to 0.09 at 400{,}000. The ATE network is more accurate at every size, with median error 0.30 to 0.41 of the CATE network's at $z=0$ (0.054 at 100{,}000 pairs), so a conditional effect needs several times the simulations of an average one for the same accuracy.

\begin{figure}[ht]
\centering
\includegraphics[width=\linewidth,trim=0bp 0bp 0bp 362bp,clip]{gbc_ate.pdf}
\caption{The conditional effect. (e) Exact posterior mean and 90 percent band of $\tau(z)$ at the observed data set, with the median and the 5 and 95 percent points of the CATE network. (f) Posterior sd of $\tau(z)$ from the exact posterior, from the exact marginals combined as if independent, from triangular draws and from independent draws.}
\label{comp:fig:cate}
\end{figure}

The second decision matters for the same curve. The exact posterior correlation of $\Delta\mu=\mu_1-\mu_0$ and $\Delta\alpha=\alpha_1-\alpha_0$ is 0.368. Triangular draws, from a network for $\Delta\mu$ and one for $\Delta\alpha$ given $S$ and the drawn $\Delta\mu$, have correlation 0.390; independent draws take $\Delta\alpha$ from a network given $S$ alone. At $z=-1.5$ and 1.5 the sd of $\tau(z)$ is 0.990 and 1.017 times the exact sd for triangular draws, 1.251 and 0.862 for independent draws, and 1.247 and 0.858 for the exact marginals combined as if independent: the fault is the lost dependence, not the networks (Figure~\ref{comp:fig:cate}f). We chose the assignment intercept so that the components are correlated: at intercept 0 the correlation is 0.001.

The example is narrower than the recipe: the summary is sufficient, the prior normal, the only confounder one observed covariate, and a network serves one design. It has no AIPW summary and no cross fitted nuisances, and we have not tried the coupled horseshoe prior with a network.
